\section{The scalar rate and its curvature}
\label{scalar:section}

A soft coordinate fixes the rate used throughout the proof.  We need
its convexity, two bounds on its entropy-scaled curvature, and one
comparison of logarithmic slopes.  These facts follow from a positive
hyperbolic integral and two elementary moment comparisons.

All logarithms are natural.  For a sign with mean $r\in[-1,1]$, put
\begin{equation}
\label{scalar:entropy-definitions}
 L=\log 2,\qquad
 \psi(r)=\frac{(1+r)\log(1+r)+(1-r)\log(1-r)}2,\qquad
 H(r)=L-\psi(r),
\end{equation}
with $0\log0=0$.  Thus $H$ is binary entropy and $\psi$ is its
gap from the entropy of a fair sign.  Write
$V(r)=\operatorname{atanh}r$ for $-1<r<1$.  Then
\begin{equation}
\label{scalar:entropy-derivatives}
 \psi'(r)=V(r),\qquad \psi''(r)=\frac1{1-r^2},\qquad H'(r)=-V(r).
\end{equation}
The restriction of $\psi$ to $[0,1)$ increases bijectively onto
$[0,L)$.  Define the production rate and its surplus over the linear
rate by
\begin{equation}
\label{scalar:profile-definitions}
 F(\psi(r))=rV(r)\qquad(0\le r<1),
\end{equation}
\begin{equation}
\label{scalar:surplus-definition}
 S(I)=
 \begin{cases}
  F(I)-2I,&0\le I<L,\\
  0,&I<0.
 \end{cases}
\end{equation}

\begin{lemma}[Elementary entropy bounds]
\label{scalar:entropy-bounds}
For $|r|\le1$,
\begin{equation}
\label{scalar:entropy-quadratic-bounds}
 \frac{r^2}{2}\le\psi(r)\le Lr^2,\qquad H(r)\ge L(1-r^2).
\end{equation}
The quotient $\psi(r)/r^2$ is strictly increasing for $0<r<1$.
Moreover,
\begin{equation}
\label{scalar:logarithm-constants}
 \frac23<L<\frac{25}{36}<\frac7{10},\qquad
 \frac12<M:=\frac{4L^2}{3}<\frac23.
\end{equation}
\end{lemma}
\begin{proof}
Integrating the geometric series for $\psi''$ twice gives
\begin{equation}
\label{scalar:entropy-series}
 \psi(r)=\sum_{n\ge1}\frac{r^{2n}}{2n(2n-1)}.
\end{equation}
The first term gives the lower bound.  The coefficients are positive
and sum to $\psi(1)=L$, so $r^{2n}\le r^2$ gives the upper bound.
The same series proves the quotient's strict increase.
Subtracting from $L$ gives the bound for $H$.
Finally, $L=2\operatorname{atanh}(1/3)$ implies
\[
 \frac23<L
 =\frac23+2\sum_{n\ge1}\frac{3^{-(2n+1)}}{2n+1}
 <\frac23+\frac23\sum_{n\ge1}3^{-(2n+1)}
 =\frac{25}{36}.
\]
Hence $16/27<M<625/972<2/3$, as required.
\end{proof}

\begin{lemma}[Convexity of the rate and surplus]
\label{scalar:profile-convexity}
The function $F$ is smooth on $[0,L)$, with right derivatives at zero,
and
\begin{equation}
\label{scalar:production-convexity}
 F(0)=0,\qquad F'(0)=2,\qquad F''(0)=\frac43,\qquad
 F''(I)\ge\frac43,\qquad F'''(I)\ge0.
\end{equation}
The zero-extended surplus $S$ is nonnegative, nondecreasing, convex, and
$C^1$ on $(-\infty,L)$, with $S(0)=S'(0)=0$.  It is not $C^2$ at zero.
\end{lemma}
\begin{proof}
Put $r=\tanh v$ and $I=\psi(r)$.  Direct differentiation gives
\[
 I=v\tanh v-\log\cosh v,\qquad
 \frac{dI}{dv}=v\operatorname{sech}^2v,\qquad
 F'(I)=1+\frac{\sinh(2v)}{2v}.
\]
For $v>0$, another differentiation yields
\begin{align}
\label{scalar:production-second-derivative}
 F''(\psi(r))
 &=\frac{(1+r^2)V(r)-r}{(1-r^2)^2V(r)^3}\\
 &=\cosh^2v\,A(v),\qquad
 A(v):=2\int_0^1(1-s^2)\cosh(2vs)\,ds.
\label{scalar:production-positive-series}
\end{align}
Indeed, integration by parts evaluates the integral as
$[2v\cosh(2v)-\sinh(2v)]/(2v^3)$.
The integral gives $A(v)\ge A(0)=4/3$ and $A'(v)>0$ for $v>0$.
Thus $F''\ge4/3$ and $F'''>0$ for $I>0$.
Both $A(v)$ and $I$ are analytic in $v^2$, and
$I=v^2/2+O(v^4)$, so the inverse function theorem gives smooth
extension to $I=0$.  Continuity gives $F'''(0)\ge0$.
The identities $S(0)=S'(0)=0$ and $S''=F''>0$ for $I>0$
prove the assertions about its zero extension.  Its second derivative
jumps from $0$ to $4/3$ at zero.
\end{proof}

\begin{lemma}[Growth of the curvature]
\label{scalar:curvature-growth}
The function $I\mapsto e^{-2I}F''(I)$ is strictly increasing on
$[0,L)$.  For $I>0$ it satisfies
\[
 \frac{F'''(I)}{F''(I)}>2.
\]
\end{lemma}
\begin{proof}
Using \eqref{scalar:production-positive-series}, for $v>0$,
\[
 \frac{d}{dI}\log F''(I)
 =\frac{\sinh(2v)}v
  +\frac{\cosh^2v}{v}\frac{A'(v)}{A(v)}>2.
\]
Here $A'(v)\ge0$ and $\sinh(2v)>2v$.  Continuity at zero
extends the asserted strict increase to $[0,L)$.
\end{proof}

\begin{lemma}[Scaled entropy curvature]
\label{scalar:scaled-curvature}
For $0<h\le L$, put
\begin{equation}
\label{scalar:scaled-curvature-definition}
 P(h)=h^2F''(L-h).
\end{equation}
Then
\begin{equation}
\label{scalar:scaled-curvature-range}
 \frac12\le P(h)\le M=\frac{4L^2}{3},\qquad P(L)=M.
\end{equation}
\end{lemma}
\begin{proof}
Write $h=H(r)$, $t=r^2$, and $x=V(r)/r$, with $x=1$ at
$r=0$.  Elementary integration gives
\begin{equation}
\label{scalar:entropy-moment-integrals}
 x=\int_0^1\frac{du}{1-tu^2},\qquad
 \frac{H(r)}{1-t}=\int_0^1\frac{du}{(1+u)(1-tu^2)}.
\end{equation}
For the second identity, use the partial fraction identity
\[
 \frac1{(1+u)(1-tu^2)}
 =\frac1{1-t}\left(\frac1{1+u}+\frac{t(u-1)}{1-tu^2}\right)
\]
and $H(r)=L-\frac12\log(1-r^2)-rV(r)$.
Let $\mu_t$ have density $[x(1-tu^2)]^{-1}$ on $[0,1]$, and set
\[
 b_t=\E_{\mu_t}\frac1{1+u},\qquad n_t=\E_{\mu_t}u^2.
\]
Then $b_t=H(r)/[(1-t)x]$ and
$n_t=(V(r)-r)/(r^2V(r))$, with continuous values at zero.
Substitution into \eqref{scalar:production-second-derivative} gives
\[
 P(H(r))=b_t^2(1+n_t).
\]
Jensen's inequality and Cauchy--Schwarz imply
\[
 b_t\ge\frac1{1+\E_{\mu_t}u}\ge\frac1{1+\sqrt{n_t}},
 \qquad
 P(H(r))\ge\frac{1+n_t}{(1+\sqrt{n_t})^2}\ge\frac12.
\]

For the upper bound, a density proportional to $w(u)/(1-tu^2)$,
with $w$ independent of $t$, satisfies
\begin{equation}
\label{scalar:covariance-differentiation}
 \frac{d}{dt}\E_t f(u)
 =\operatorname{Cov}_t\left(f(u),\frac{u^2}{1-tu^2}\right).
\end{equation}
Oppositely monotone functions have nonpositive covariance, since
for independent copies $U,U'$,
\[
 2\operatorname{Cov}(f(U),g(U))
 =\E[(f(U)-f(U'))(g(U)-g(U'))].
\]
Applying this with $w=1$ shows $1/2\le b_t\le b_0=L$.
The same identity shows that $n_t$ increases, so we must control
its growth together with the decrease of $b_t$.  Compare
$b_t-1/2$ with $1-n_t$ by reweighting $\mu_t$ by $1-u^2$
to obtain $\nu_t$.  The identity
\[
 \beta_t:=\frac{b_t-1/2}{1-n_t}
 =\frac12\E_{\nu_t}(1+u)^{-2},\qquad
 d\nu_t(u)\ \propto\ \frac{1-u^2}{1-tu^2}\,du
\]
and \eqref{scalar:covariance-differentiation} now give
$0<\beta_t\le\beta_0=\frac32(L-\frac12)$.
Consequently
\[
 1+n_t\le2-\frac{b_t-1/2}{\beta_0},
\]
and hence
\[
 P(H(r))\le f(b_t),\qquad
 f(b)=b^2\left[2-\frac{b-1/2}{\beta_0}\right].
\]
For $1/2\le b\le L$,
\[
 f'(b)=\frac b{\beta_0}(4\beta_0+1-3b)
 \ge\frac b{\beta_0}(3L-2)>0.
\]
Thus $P(H(r))\le f(L)=4L^2/3$.
At $r=0$, the uniform law has $b_0=L$ and $n_0=1/3$, which
also gives $P(L)=4L^2/3$.
\end{proof}
